\section{Conclusão}

As diferenças entre produção e demanda se compõem de relações com volumes, concentrações e contrapartes distintos. No cenário analisado, energia distribui emissões entre diversos produtos finais e construção reúne atribuições de outras origens, mantendo grande parcela própria na diagonal. A comparação com os requerimentos como fornecedor mostra por que magnitude ambiental e posição produtiva precisam ser distinguidas. O backbone facilita localizar relações salientes, desde que sua cobertura acompanhe a interpretação.

A contribuição é um procedimento auditável para passar dos totais setoriais às relações que os conectam. Ele orienta a investigação conjunta de origens e produtos finais para mitigação, sem escolher instrumentos ou estimar abatimento. A aproximação das intensidades limita as conclusões ambientais; a conservação contábil e a seleção de arestas não eliminam essa condição.
